% !TeX root = ../papers/nc-p3.tex
\section{階数一切断の長さと非被約な例外鉛筆}
\label{sec:finite-marked-secant-pencil}
前節の共通行列四平面を $E$ とし、$\Sigma=\mathbb P(E)\cap(\mathbb P^2\times\mathbb P^2)$ を
階数1のスキーム切断とする。接線を数えるためスキームの長さを使う。
\begin{theorem}\label{thm:finite-marked-secant-pencil}
$\Sigma$ が有限なら長さは4以下である。長さ2以下のとき、
印の接次元が0または1なら、他の局所成分への仮定なしに前節の $T$ は反対対称性を与える。
印の接次元が2でも、その非零接方向に前節の一次元条件が成立すれば同じである。
従って両側でこの判定を通れば全次数AS判定ができる。
\end{theorem}
\begin{lemma}\label{lem:finite-adjacent-word-support}
長さ20の有限辺で $C=U_0U_1\subset D$ とし、$0\ne f\in C$ の支持を
$Z=\operatorname{Spec}(D/\operatorname{ann} f)$ とする。
$Z$ 上での二つの座標系の線型包の次元を $a,b$（射影次元に1を加えた値）とすると
\[
 a+b\le d=\operatorname{length} Z\le\binom{a+b-2}{a-1}.
\]
特に $a,b\ge3$, $d\ge6$ で、$d=6$ なら $a=b=3$、$Z$ は二平面の積での
四つの双一次式の完全交叉である。第二位置のどの直線 $L$ に対しても
$d-\operatorname{length}(Z\cap L)\ge b$ が成り立つ。ここで交叉は第二射影で取る。
\end{lemma}
\begin{proof}
$\tau(f\,\cdot)$ は商の完全跡であり、$C$ が $C$ と跡直交するので
二座標空間の制限像はこの対合で互いに直交する。
両射影の包 $\mathbb P^{a-1}\times\mathbb P^{b-1}$ に元の六式を制限すれば
基点スキームは有限で $Z$ を含む。従ってそのうち $a+b-2$ 本の一般線型結合が
固有に交わる（有限基点の外で各段の正次元成分を避ける）。
交点数は表示の二項係数である。$a$ または $b$ が1,2なら二つの不等式は矛盾する。
$d=6$ では全ての不等式が等号となり、完全交叉に含まれる同じ長さの部分は全体である。
$Z\cap L$ は第一射影の包と $L$ の積で同様に長さ $a$ 以下なので最後の主張を得る。
\end{proof}
\begin{proof}[階数一切断と固定鉛筆]
$\mathbb P(E)$ を定める五超平面のうち一般の四線型結合でSegre四次元多様体を切ると、
有限基点仮定により長さ6の完全交叉 $Y$ になる。これは切断の射影四平面を張る。
実際 $(1,1)$ で消える切断は四つの定義式で尽きる（Koszul分解で $j=2,3$ の
$\mathcal O(1-j,1-j)$ は全コホモロジーが零、$j=4$ の非零コホモロジーは次数4だけである）。
$Y$ での $(1,1)$ 切断の像は五次元で、随伴 $\omega_Y=\mathcal O_Y(1,1)$ の完全跡の
核である。これは同じKoszul分解の余核への写像が双対で単位の制限になることによる。
残った超平面の切断は非零で、生成するイデアルの長さが1なら、ある局所因子のsocleを張る。
完全跡はこの元で非零なので矛盾する。よって残余長さは2以上、$\Sigma$ の長さは4以下である。

左右の指定核 $l,r$ を消す二次元鉛筆は二行二列の行列表示を持つ。その射影直線を
階数1で切る二次式は恒等零ではない（$\Sigma$ は有限）。従って長さ2の部分を切り出す。
$\Sigma$ の長さが2以下なら、そのような直線がある場合は長さが2で、直線は
$\Lambda=\langle\Sigma\rangle$ に一意である。二行二列の一般元は階数2なので
共通左右核も一意である。この一意性は局所Artin基底上でも成立する。
実際、二つの元が剰余体で独立で、共通左右核が直和因子の直線なら、
二次式による切断は相対射影直線で有限平坦な次数2の因子となる。
局所的には非根を無限遠に選び、単元を先頭係数とする二次式の商で確認できる。
この因子は $\Sigma$ の基底変換に含まれ、長さ2の自由加群間の全射は同型である。
従って相対鉛筆も相対核も一定である。$\Sigma$ が非被約長さ2でも、その一次式の核が
張る射影直線は一意なので同じ証明でよい。

以下、例外鉛筆があるとし、その左核が印からの射影で指定する位置1の直線を $L$ とする。
印以外の点で前節の $A,B$ の比例からのずれを調べる。
縮約核が二次元以上ならちょうど二次元である。非零の左縮約は階数2以上であることが
有限辺での純関係の不在から従う（定理\ref{thm:finite-cubic-separators}の証明）。
$\lambda(A)$ を単元として三次元の $\ker\lambda$ にある差 $Z$ を取ると、
その六本の縮約条件の階数は剰余体上で2である。二つの主元で消去した後は
$Z=f_0 V$、$V$ の一成分は1、$V$ と $A$ は剰余体で独立となる。
$f_0\ne0$ なら $\operatorname{ann}(f_0)$ を法として $A,V$ は実際に二本の共通核を消す。
相対鉛筆の一意性よりこの商で $x_1$ は $L$ にある。ゆえに $I_L Z=0$。
縮約核が一次元のときは最初から $Z=0$ である。

印が曲線状なら局所環を $k[t]/(t^m)$ とする。印では $A,B$ は比例するので
$Z\ne0$ なら $Z=t^s V$, $s\ge1$ と書け、$V$ は $A$ と剰余体で独立である。
射影三座標 $l,r$ はいずれも $t$ で割れ、それぞれの商は一成分が単元である。
$m-s\ge2$ なら $t^{m-s-1}$ を法として、$A,V$ は $l/t,r/t$ を共通核に持つ。
よって $l/t$ はその商で固定方向に入り、$l$ は $t^{m-s}$ を法としてその方向に入る。
$m-s=1$ のときも $l$ の剰余は零である。結局 $I_L Z=0$ が印でも成立する。
印が被約なら差は零である。従って各第二辺関係を $(x_1,T x_0)$ で評価した $f\in C$ は
全て $I_L f=0$ となる。非零なら補題の $b\ge3$ に反する。

印が接次元2で全接方向の条件が成立する場合、前節の行加群の証明から印で
$\mathfrak m^3 f=0$ である。$f\ne0$ と仮定し、補題の支持 $Z$ を考える。
印以外は $L$ 上にあり、印の局所成分があればGorensteinで極大イデアルの三乗が零である。
従ってその長さは4以下である。実際二乗が非零ならsocleそのもので次元1である。
$Z$ の印の接次元が1以下ならここで $L$ を課して失う長さは高々2であり、
補題の $d-\operatorname{length}(Z\cap L)\ge b\ge3$ に反する。
接次元2なら長さ4で、その接平面は元の $X$ の接平面である。
$L$ がこの平面に含まれるなら $I_L$ の一次部分の像は高々一次元で、
二乗部分も一次元なので同じ矛盾を得る。
$L$ がこの平面に含まれなければ、失う長さは高々3で、補題から $b=3$。
従って第二射影の包はまさに印の接平面である。
$d\ge6$ なので印以外にも支持点があり、それと印を結ぶ $L$ はその平面に入らねばならず矛盾する。
以上で評価は全て零となり、前節と同じグラフの議論で結論を得る。
\end{proof}
